\clearpage
\section{Design}

The design separates \textbf{game semantics} from \textbf{execution strategy}.
All versions use the same game model, board, view, and physics rules. The
sequential version is the reference. The two concurrent versions add an
asynchronous runtime and parallel physics without changing the domain rules.

\subsection{Overall architecture}

The concurrent architecture follows one main control path:

\begin{center}
\begin{tikzpicture}[
    scale=0.86,
    transform shape,
    class/.style={draw, rounded corners=2pt, align=center, minimum height=8mm,
                  minimum width=25mm, inner xsep=5pt, font=\small},
    presentation/.style={class, fill=green!10, draw=green!45!black},
    runtime/.style={class, fill=orange!12, draw=orange!55!black},
    domainclass/.style={class, fill=blue!8, draw=blue!50!black},
    physics/.style={class, fill=yellow!15, draw=yellow!50!black},
    interface/.style={class, fill=white, dashed, line width=0.7pt},
    uses/.style={-{Latex[length=2mm]}, thick},
    implements/.style={-{Latex[length=2mm]}, dashed, thick},
    node distance=7mm and 8mm
]
    \node[presentation] (viewFrame) {ViewFrame};
    \node[presentation, left=of viewFrame] (view) {View};
    \node[presentation, right=of viewFrame] (viewModel) {ViewModel};

    \node[interface, below=10mm of viewFrame] (gameRuntime) {GameRuntime\\\scriptsize interface};
    \node[runtime, below left=7mm and 5mm of gameRuntime] (threadedRunner) {Threaded\\GameRunner};
    \node[runtime, below right=7mm and 5mm of gameRuntime] (taskRunner) {TaskBased\\GameRunner};
    \node[runtime, below=22mm of gameRuntime] (gameLoop) {GameLoop};
    \node[runtime, left=of gameLoop] (botAgent) {BotAgent};

    \node[domainclass, below left=11mm and 5mm of gameLoop] (gameModel) {GameModel};
    \node[domainclass, right=of gameModel] (board) {Board};
    \node[interface, below=9mm of board] (physicsStepper) {PhysicsStepper\\\scriptsize interface};

    \node[physics, below=9mm of physicsStepper] (threadedEngine) {Threaded\\PhysicsEngine};
    \node[physics, left=5mm of threadedEngine] (sequentialEngine) {Sequential\\PhysicsEngine};
    \node[physics, right=5mm of threadedEngine] (taskEngine) {TaskBased\\PhysicsEngine};

    \node[left=7mm of view, font=\bfseries\scriptsize] {Presentation};
    \node[left=7mm of threadedRunner, font=\bfseries\scriptsize] {Runtime};
    \node[left=7mm of gameModel, font=\bfseries\scriptsize] {Domain};
    \node[left=7mm of sequentialEngine, font=\bfseries\scriptsize] {Physics};

    \draw[uses] (view) -- node[above, font=\scriptsize] {creates} (viewFrame);
    \draw[uses] (viewModel) -- node[above, font=\scriptsize] {state} (viewFrame);
    \draw[uses] (viewFrame) -- node[left, font=\scriptsize] {callbacks} (gameRuntime);
    \draw[implements] (threadedRunner) -- (gameRuntime);
    \draw[implements] (taskRunner) -- (gameRuntime);
    \draw[uses] (threadedRunner) -- (gameLoop);
    \draw[uses] (taskRunner) -- (gameLoop);
    \draw[uses] (botAgent) -- node[above, font=\scriptsize] {send commands} (gameLoop);
    \draw[uses] (gameLoop) -- (gameModel);
    \draw[uses] (gameModel) -- (board);
    \draw[uses] (board) -- node[right, font=\scriptsize] {delegates} (physicsStepper);
    \draw[implements] (sequentialEngine) -- (physicsStepper);
    \draw[implements] (threadedEngine) -- (physicsStepper);
    \draw[implements] (taskEngine) -- (physicsStepper);
    \draw[uses] (gameLoop.east) -| ([xshift=8mm]viewModel.east)
        |- node[pos=0.75, right, font=\scriptsize] {snapshots} (viewModel.east);

    \node[below=3mm of threadedEngine, font=\scriptsize] {
        \tikz\draw[uses] (0,0) -- (5mm,0); uses / data flow
        \qquad
        \tikz\draw[implements] (0,0) -- (5mm,0); implements
    };
\end{tikzpicture}
\end{center}

The launcher assembles these components. User actions become callbacks to the
selected \texttt{GameRuntime}. The runtime delegates them to
\texttt{GameLoop}, which orders commands, advances \texttt{GameModel}, and
publishes the completed state. The model owns match rules, while
\texttt{Board} owns physical state. The \texttt{PhysicsStepper} interface
selects the sequential, threaded, or task-based physics engine.

The output follows the opposite direction. A completed game state is copied
into a \texttt{RuntimeGameSnapshot}, then into \texttt{ViewModel}. The view
therefore draws copied data and never accesses mutable balls while physics is
running.

\subsection{Single-writer runtime and command flow}

The central ownership rule is: \textbf{GameModel has one writer}. In the
concurrent versions, only the controller executing \texttt{GameLoop.tick(...)}
may advance the match. Human input and the bot do not receive direct write
access. They submit operations to \texttt{CommandMailbox}.

The mailbox is a FIFO monitor. Each submitted command carries a receipt that
is completed after the command is executed or rejected. At the start of a
tick, the controller drains the queue. It then advances the game and publishes
one new snapshot:

\begin{center}
\textbf{drain commands} $\rightarrow$ \textbf{advance model}
$\rightarrow$ \textbf{publish snapshot}
\end{center}

This order gives each accepted input a clear position between simulation
steps. It also prevents the view and bot from observing an intermediate state.

\begin{center}
\begin{tabular}{|>{\bfseries}p{3.1cm}|p{5.2cm}|p{5.2cm}|}
    \hline
    State & Writer & Readers or publication boundary \\
    \hline
    Game rules and scores & GameLoop controller through GameModel & RuntimeGameSnapshot \\
    \hline
    Balls and holes & Current physics step & Board snapshots after the step \\
    \hline
    Pending commands & CommandMailbox monitor & GameLoop during queue draining \\
    \hline
    Rendering state & Launcher through ViewModel & ViewFrame \\
    \hline
\end{tabular}
\end{center}

\subsection{Execution strategies}

The concurrent versions preserve the same control protocol but use different
mechanisms to execute it.

\begin{center}
\begin{tabular}{|>{\bfseries}p{3.1cm}|p{5.2cm}|p{5.2cm}|}
    \hline
    Aspect & Platform-thread version & Task-based version \\
    \hline
    Controller & Dedicated platform thread & Single-thread ScheduledExecutorService \\
    \hline
    Bot & Dedicated platform thread & Task on a dedicated ExecutorService \\
    \hline
    Physics work & Long-lived PhysicsWorker threads & Tasks submitted to a fixed thread pool \\
    \hline
    Phase completion & WorkerCompletionMonitor & Future.get() barriers \\
    \hline
    Resource release & Workers are interrupted and joined & Executors are shut down and awaited \\
    \hline
\end{tabular}
\end{center}

The two physics engines intentionally remain separate. This keeps the required
programming models explicit: \texttt{ThreadedPhysicsEngine} directly owns and
coordinates platform threads, while \texttt{TaskBasedPhysicsEngine} expresses
the same work as executor tasks. Their phase structure and numerical rules are
kept equivalent so performance comparisons remain meaningful.

\paragraph{Lifecycle and shutdown.}
Runners own their active resources. Shutdown reverses startup: \textbf{stop
ticks, close the mailbox, publish the final snapshot, stop bot and controller,
then close the physics engine}. This rejects new work and releases pending
command receipts and snapshot waiters.

\clearpage
\subsection{Phased physics design}

One call to \texttt{step(...)} holds the board as a single mutable unit. A long
elapsed interval is divided into fixed sub-steps. Each sub-step follows ordered
phases. The following parametric P/T net distinguishes sequential transitions
from parallel fork--join phases without modelling scheduler internals.

\begin{center}
\begin{tikzpicture}[
scale=1.02,
transform shape,
place/.style={circle, draw=black!70, minimum size=5.5mm, inner sep=0pt},
transition/.style={rectangle, draw=black!75, fill=black!70,
minimum width=2.2mm, minimum height=8mm, inner sep=0pt},
worker/.style={transition, fill=blue!55!black, draw=blue!70!black},
serial/.style={transition, fill=orange!75!black, draw=orange!80!black},
flow/.style={-{Latex[length=1.6mm]}, line width=0.6pt},
phase/.style={font=\tiny, align=center},
state/.style={font=\tiny, align=center, text width=18mm},
wlabel/.style={font=\tiny}
]
% Sequential control backbone, followed by representative parallel sub-nets.
\node[place, tokens=1] (ready) at (0,0) {};
\node[transition] (moveFork) at (1.2,0) {};
\node[place] (moveInOne) at (2.4,0.5) {};
\node[place] (moveInLast) at (2.4,-0.5) {};
\node[worker] (moveOne) at (3.6,0.5) {};
\node[worker] (moveLast) at (3.6,-0.5) {};
\node[place] (moveOutOne) at (4.8,0.5) {};
\node[place] (moveOutLast) at (4.8,-0.5) {};
\node[transition] (moveJoin) at (6,0) {};
\node[place] (positionsUpdated) at (7.2,0) {};
\node[serial] (holes) at (8.6,0) {};
\node[place] (holesProcessed) at (10,0) {};

\draw[flow] (ready) -- (moveFork);
\draw[flow] (moveFork) -- (moveInOne);
\draw[flow] (moveInOne) -- (moveOne);
\draw[flow] (moveOne) -- (moveOutOne);
\draw[flow] (moveOutOne) -- (moveJoin);
\draw[flow] (moveFork) -- (moveInLast);
\draw[flow] (moveInLast) -- (moveLast);
\draw[flow] (moveLast) -- (moveOutLast);
\draw[flow] (moveOutLast) -- (moveJoin);
\draw[flow] (moveJoin) -- (positionsUpdated);
\draw[flow] (positionsUpdated) -- (holes);
\draw[flow] (holes) -- (holesProcessed);

\node[state, below=1.5mm of ready] {tick ready};
\node[phase] at (3.6,0) {$\vdots$};
\node[wlabel, above left=0.5mm and 1.5mm of moveOne] {$W_1$};
\node[wlabel, below left=0.5mm and 1.5mm of moveLast] {$W_R$};
\node[phase, above=7mm of moveOne] {Update disjoint active-ball ranges in parallel};
\node[state, below=1.5mm of positionsUpdated] {positions updated};
\node[phase, above=1.5mm of holes] {Resolve hole interactions\\sequentially};
\node[state, below=1.5mm of holesProcessed] {holes resolved};

% Parallel private-grid construction, followed by deterministic merge.
\node[transition] (gridFork) at (8.8,-3.4) {};
\node[place] (gridInOne) at (7.6,-2.9) {};
\node[place] (gridInLast) at (7.6,-3.9) {};
\node[worker] (gridOne) at (6.4,-2.9) {};
\node[worker] (gridLast) at (6.4,-3.9) {};
\node[place] (gridOutOne) at (5.2,-2.9) {};
\node[place] (gridOutLast) at (5.2,-3.9) {};
\node[transition] (gridJoin) at (4,-3.4) {};
\node[place] (privateGridsBuilt) at (2.8,-3.4) {};
\node[serial] (mergeGrid) at (1.5,-3.4) {};
\node[place] (orderedGridReady) at (0.2,-3.4) {};

\draw[flow] (holesProcessed) |- (gridFork);
\draw[flow] (gridFork) -- (gridInOne);
\draw[flow] (gridInOne) -- (gridOne);
\draw[flow] (gridOne) -- (gridOutOne);
\draw[flow] (gridOutOne) -- (gridJoin);
\draw[flow] (gridFork) -- (gridInLast);
\draw[flow] (gridInLast) -- (gridLast);
\draw[flow] (gridLast) -- (gridOutLast);
\draw[flow] (gridOutLast) -- (gridJoin);
\draw[flow] (gridJoin) -- (privateGridsBuilt);
\draw[flow] (privateGridsBuilt) -- (mergeGrid);
\draw[flow] (mergeGrid) -- (orderedGridReady);

\node[phase] at (6.4,-3.4) {$\vdots$};
\node[wlabel, above left=0.5mm and 1.5mm of gridOne] {$W_1$};
\node[wlabel, below left=0.5mm and 1.5mm of gridLast] {$W_R$};
\node[phase, above=7mm of gridOne] {Build one private spatial grid per range in parallel};
\node[state, below=1.5mm of privateGridsBuilt] {local grids built};
\node[phase, above=1.5mm of mergeGrid] {Merge grids and order\\occupied cells sequentially};
\node[state, below=1.5mm of orderedGridReady] {grid merged and ordered};

% Parallel cell inspection, followed by sequential contribution accumulation.
\node[transition] (pairFork) at (1.5,-6.8) {};
\node[place] (pairInOne) at (2.8,-6.3) {};
\node[place] (pairInLast) at (2.8,-7.3) {};
\node[worker] (pairOne) at (4,-6.3) {};
\node[worker] (pairLast) at (4,-7.3) {};
\node[place] (pairOutOne) at (5.2,-6.3) {};
\node[place] (pairOutLast) at (5.2,-7.3) {};
\node[transition] (pairJoin) at (6.4,-6.8) {};
\node[place] (pairsCollected) at (7.6,-6.8) {};
\node[serial] (computeDeltas) at (8.8,-6.8) {};
\node[place] (deltasAccumulated) at (10,-6.8) {};

\draw[flow] (orderedGridReady) |- (pairFork);
\draw[flow] (pairFork) -- (pairInOne);
\draw[flow] (pairInOne) -- (pairOne);
\draw[flow] (pairOne) -- (pairOutOne);
\draw[flow] (pairOutOne) -- (pairJoin);
\draw[flow] (pairFork) -- (pairInLast);
\draw[flow] (pairInLast) -- (pairLast);
\draw[flow] (pairLast) -- (pairOutLast);
\draw[flow] (pairOutLast) -- (pairJoin);
\draw[flow] (pairJoin) -- (pairsCollected);
\draw[flow] (pairsCollected) -- (computeDeltas);
\draw[flow] (computeDeltas) -- (deltasAccumulated);

\node[phase] at (4,-6.8) {$\vdots$};
\node[wlabel, above left=0.5mm and 1.5mm of pairOne] {$W_1$};
\node[wlabel, below left=0.5mm and 1.5mm of pairLast] {$W_R$};
\node[phase, above=7mm of pairOne] {Inspect disjoint occupied-cell ranges in parallel};
\node[state, below=1.5mm of pairsCollected] {collision pairs collected};
\node[phase, above=1.5mm of computeDeltas] {Compute and accumulate collision\\contributions sequentially};
\node[state, below=1.5mm of deltasAccumulated] {collision deltas accumulated};

% R=1 gives sequential apply; R>1 enables concurrent branches.
\node[transition] (applyFork) at (8.8,-10.2) {};
\node[place] (applyInOne) at (7.6,-9.7) {};
\node[place] (applyInLast) at (7.6,-10.7) {};
\node[worker] (applyOne) at (6.4,-9.7) {};
\node[worker] (applyLast) at (6.4,-10.7) {};
\node[place] (applyOutOne) at (5.2,-9.7) {};
\node[place] (applyOutLast) at (5.2,-10.7) {};
\node[transition] (applyJoin) at (4,-10.2) {};
\node[place] (complete) at (2.8,-10.2) {};

\draw[flow] (deltasAccumulated) |- (applyFork);
\draw[flow] (applyFork) -- (applyInOne);
\draw[flow] (applyInOne) -- (applyOne);
\draw[flow] (applyOne) -- (applyOutOne);
\draw[flow] (applyOutOne) -- (applyJoin);
\draw[flow] (applyFork) -- (applyInLast);
\draw[flow] (applyInLast) -- (applyLast);
\draw[flow] (applyLast) -- (applyOutLast);
\draw[flow] (applyOutLast) -- (applyJoin);
\draw[flow] (applyJoin) -- (complete);

\node[phase] at (6.4,-10.2) {$\vdots$};
\node[wlabel, above left=0.5mm and 1.5mm of applyOne] {$W_1$};
\node[wlabel, below left=0.5mm and 1.5mm of applyLast] {$W_R$};
\node[phase, above=7mm of applyOne] {Apply final deltas: $R=1$ sequential, $R>1$ parallel};
\node[state, below=1.5mm of complete] {complete};

\end{tikzpicture}

\small \textcolor{blue!55!black}{Blue transitions} are concurrently enabled
range operations. \textcolor{orange!75!black}{Orange transitions} are
sequential coordinator operations. Dark transitions dispatch and join ranges.
\end{center}

Each fork represents $R$ independent range instances. The multiplicity appears
on the arcs, while places still encode the current marking and transitions
remain the only atomic state changes. The corresponding blue transitions can
fire in any order or concurrently. Private input places ensure that each range
fires exactly once; the join requires every completion token. The net is
therefore 1-safe and prevents the next phase from starting early. For final
delta application, $R=1$ models the sequential fast path, while $R>1$ models
the parallel path selected for large touched-ball sets.

The spatial grid is the broad phase. It inspects only the cell containing a
ball and selected neighbours, avoiding every possible pair on sparse boards.
Private grids and pair lists also remove parallel write contention.

\subsection{Determinism and serialized commit}

Parallel workers never update scores or remove balls. They produce private
lists of confirmed pairs, which the coordinator merges and sorts before
calculating effects. The order is independent of map iteration and completion
timing.

Several collisions can affect the same ball. Immediate application would cause
competing writes and schedule-dependent floating-point results. The engines
instead accumulate position and velocity corrections before applying them:

\begin{center}
\textbf{parallel computation, deterministic ordering, serialized commit}
\end{center}

The final delta application may become parallel again because each touched
ball has one accumulated result and can be updated independently. For small
sets it remains sequential, avoiding a barrier whose cost would exceed the
available work.

\subsection{View boundary}

\texttt{RuntimeGameSnapshot} publishes copied game and physics state;
\texttt{ViewModel} converts it into synchronized rendering records, so the view
never reads the model directly. Keyboard and mouse callbacks produce a shot
vector, but the runtime validates it on the model owner, not in
\texttt{ViewFrame}. Frame identifiers let the launcher await each repaint, so
the UI displays complete states in order.
